\subsection{old}




\begin{lstlisting}[language=C, caption={Netduino Plus 2 machine initialization}]
static void netduinoplus2_init(MachineState *machine)
{
    DeviceState *dev;
    Clock *sysclk;

    /* This clock doesn't need migration because it is fixed-frequency */
    sysclk = clock_new(OBJECT(machine), "SYSCLK");
    clock_set_hz(sysclk, SYSCLK_FRQ);

    /* Instantiate STM32F405 SoC */
    dev = qdev_new(TYPE_STM32F405_SOC);
    object_property_add_child(OBJECT(machine), "soc", OBJECT(dev));
    qdev_connect_clock_in(dev, "sysclk", sysclk);
    sysbus_realize_and_unref(SYS_BUS_DEVICE(dev), &error_fatal);

    /* Initialize LIS3DH on I2C1 */
    STM32F405State *s = STM32F405_SOC(dev);    
    DeviceState *lis3dh = qdev_new(TYPE_LIS3DH); // Creating LIS3DH accelerometer
    
    i2c_slave_set_address(I2C_SLAVE(lis3dh), 0x18); // Setting I2C address to 0x18    
    i2c_slave_realize_and_unref(I2C_SLAVE(lis3dh), 
                    s->i2c[0].bus, &error_fatal);    // Connecting to I2C1 bus

    qdev_connect_gpio_out(lis3dh, 0, qdev_get_gpio_in(DEVICE(&s->exti), 1));

    /* Load firmware into flash memory */
    armv7m_load_kernel(ARM_CPU(first_cpu),
                       machine->kernel_filename,
                       0, FLASH_SIZE);
}
\end{lstlisting}


%%%%%


When the LIS3DH asserts its INT1 output in response to a data-ready event, the signal propagates through the device hierarchy as follows:

\begin{enumerate}
\item \textbf{LIS3DH Assertion}: The `lis3dh\_update\_drdy\_interrupt()` function calls `qemu\_irq\_raise(s->int1)` when new acceleration data becomes available and the I1\_ZYXDA bit in CTRL\_REG3 is set.

\item \textbf{EXTI Reception}: The QEMU IRQ infrastructure routes the signal to the EXTI device's `gpio\_in[1]` input. The EXTI device updates its pending register (EXTI\_PR) and, if the corresponding interrupt mask bit (EXTI\_IMR) is enabled, propagates the interrupt to its output.

\item \textbf{NVIC Activation}: The EXTI output triggers NVIC IRQ 7 (EXTI1\_IRQn). The NVIC evaluates interrupt priority, preemption settings, and masking, then signals the ARM Cortex-M4 core to execute the exception handler.

\item \textbf{Firmware ISR Execution}: The Cortex-M4 core vectors to the EXTI1\_IRQHandler, which executes firmware-defined interrupt service logic. Typically, the ISR reads the LIS3DH's INT1\_SRC register to identify the interrupt source and clears the EXTI pending bit by writing to EXTI\_PR.

\item \textbf{Interrupt Clearance}: If the LIS3DH is configured in latched interrupt mode (LIR\_INT1 = 1 in CTRL\_REG5), the interrupt remains asserted until firmware reads INT1\_SRC. Reading this register executes:


This lowers the INT1 line, clearing the EXTI pending bit and completing the interrupt cycle.
\end{enumerate}



%%%%%%%%%%%%%%%%%%%%5


This code demonstrates several key integration patterns:

\begin{itemize}
\item \textbf{Hierarchical Device Instantiation:} The SoC is created first and realized, which triggers realization of all child I2C controllers and creation of I2C buses.

\item \textbf{Bus Reference Usage:} The \texttt{s->i2c[0].bus} pointer (populated during SoC realization) is used to attach the LIS3DH to I2C1.

\item \textbf{GPIO Routing:} The LIS3DH's INT1 output (GPIO index 0) is connected to EXTI line 1, enabling interrupt-driven firmware operation.

\item \textbf{I2C Address Configuration:} The accelerometer's 7-bit I2C address (\texttt{0x18}) is set before realization, matching the hardware configuration with SA0 pin grounded.
\end{itemize}

%%%%%%%%%%%%%%%%%%%%%%%%%

\subsection{QEMU Device Hierarchy}

The Netduino Plus 2 machine instantiates a hierarchical device tree where the LIS3DH exists as an I2C slave device attached to the STM32F405 SoC's I2C1 bus. Figure~\ref{fig:qemu_hierarchy} illustrates the complete device hierarchy:

\begin{figure}[H]
    \centering
    \begin{verbatim}
    netduinoplus2 (Machine)
    \---- STM32F405_SOC
        |---- i2c[0] (I2C1 Bus)
        |   \---- lis3dh (I2C Slave @ 0x18)
        |       |---- int1 (qemu_irq out)  ----\
        |       \---- int2 (qemu_irq out)      |
        |                                      |
        \---- exti (STM32F4XX_EXTI)            |
            |---- gpio_in[0] ← (EXTI0)         |
            |---- gpio_in[1] ← (EXTI1) <-------/
            |---- gpio_in[2] ← (EXTI2)
            \---- ...
            \---- irq_out[7] → NVIC IRQ 7 (EXTI1_IRQn)
    \end{verbatim}
    \caption{qemu hierarchy}
    \label{fig:qemu_hierarchy}
\end{figure}

This hierarchy demonstrates the signal flow from the LIS3DH interrupt outputs through the EXTI (External Interrupt) controller to the NVIC (Nested Vectored Interrupt Controller), which ultimately triggers the ARM Cortex-M4 core's interrupt service routine.





This hierarchy reflects the physical hardware organization and enables QOM-based introspection and property manipulation via the QEMU monitor.


%==== Code Quality and Testing Infrastructure ====%
\section{Code Quality and Testing Infrastructure}

\label{sec:impl_quality}

Throughout the implementation process, adherence to QEMU coding standards and validation through systematic testing were prioritized to ensure reliability and maintainability.

\subsection{Coding Standards Compliance}

All implemented code follows QEMU's coding style guidelines:

\begin{itemize}
\item \textbf{Formatting:} 4-space indentation, 80-character line limit, K\&R brace style.
\item \textbf{Naming Conventions:} Lowercase with underscores for functions and variables, capitalized for type names, prefixes indicating subsystem (\texttt{stm32f4xx\_}, \texttt{lis3dh\_}).
\item \textbf{Documentation:} Doxygen-style comments for all public functions and complex algorithms.
\item \textbf{Error Handling:} Consistent use of \texttt{Error **errp} parameters for error propagation.
\end{itemize}

Code quality was verified using QEMU's \texttt{scripts/checkpatch.pl} tool:

\begin{verbatim}
$ ./scripts/checkpatch.pl --no-tree -f hw/i2c/stm32f4xx_i2c.c
total: 0 errors, 0 warnings, 1247 lines checked
\end{verbatim}

\subsection{Testing Approach}

Validation of the implementation employed a multi-layered testing strategy:

\begin{enumerate}
\item \textbf{Unit Testing:} Individual register operations were tested using QEMU's \texttt{qtest} framework, verifying correct read/write behavior and flag updates.

\item \textbf{Firmware Integration:} Custom firmware based on STM32 HAL libraries exercised complete I2C transaction sequences (initialization, write, read, multi-byte transfers).

\item \textbf{Interrupt Verification:} Interrupt generation timing and flag clearing sequences were validated by monitoring NVIC interrupt pending registers during firmware execution.

\item \textbf{Protocol Compliance:} I2C bus signals were logged using QEMU's trace infrastructure and compared against expected I2C protocol timings from the datasheet.
\end{enumerate}

The firmware testing harness included:
\begin{itemize}
\item LIS3DH initialization sequence (configure ODR, FS, operating mode)
\item Single-byte and multi-byte register read/write operations
\item Auto-increment addressing mode validation
\item Acceleration data reading and conversion verification
\item Interrupt-driven data-ready operation
\end{itemize}

All test scenarios executed successfully, confirming functional correctness of the implemented device models.


%==== Build Artifacts and Deliverables ====%
\section{Build Artifacts and Deliverables}

\label{sec:impl_artifacts}

The implementation deliverables include:

\begin{itemize}
\item \textbf{Source Files:} 9 C source files and 9 header files totaling approximately 3,500 lines of code.
\item \textbf{Documentation:} Inline code comments, function documentation, and this implementation chapter.
\item \textbf{Build Configuration:} Modified \texttt{meson.build} files for compilation integration.
\item \textbf{Test Firmware:} STM32 HAL-based validation programs demonstrating I2C and accelerometer operation.
\end{itemize}

The complete implementation is compatible with QEMU version 9.0.0 and can be built and executed using the standard QEMU build process documented in Section~\ref{sec:impl_environment}.


%==== Summary ====%
\section{Summary}

This chapter presented the complete implementation of the STM32F405 I2C emulation and LIS3DH accelerometer device models within the QEMU framework. The implementation faithfully translates the design decisions documented in Chapter~\ref{chap:design} into working C code that adheres to QEMU's architectural patterns and coding standards.

Key implementation achievements include:

\begin{itemize}
\item Extension of the STM32F405 SoC model to instantiate three I2C controllers with correct memory mapping and interrupt routing.
\item Development of a complete I2C controller emulation featuring a finite state machine that translates firmware register operations into I2C bus events.
\item Creation of a LIS3DH accelerometer model implementing the full register map, I2C slave protocol, configuration logic, and acceleration data conversion pipeline.
\item Integration of all components within the Netduino Plus 2 machine definition, enabling end-to-end firmware testing.
\end{itemize}

The implementation has been validated through comprehensive firmware testing, demonstrating functional correctness and protocol compliance. This work establishes a foundation for future extensions, including additional I2C peripherals and enhanced testing frameworks, which will be discussed in the conclusion and future work chapter.
